\documentclass{ximera}
\input{../../../preamble.tex}


\definecolor{fvdnavy}{HTML}{071D43}
\definecolor{fvdcyan}{HTML}{20BFC6}
\definecolor{fvdcyanlight}{HTML}{EFFCFB}
\definecolor{fvdmagenta}{HTML}{F10D69}
\definecolor{fvdmagentalight}{HTML}{FFF1F7}
\definecolor{fvdtext}{HTML}{163044}





%=================================================
% Pedagogische omgevingen
% PDF: tcolorbox
% HTML: learning-box
%=================================================

\newenvironment{voorkennis}
{%
    \ifdefined\HCode
        \HCode{%
<section class="learning-box learning-box--prior">
<div class="learning-box__header">
<span class="learning-box__icon" aria-hidden="true"></span>
<span class="learning-box__title">Voorkennis</span>
</div>
<div class="learning-box__content">
        }%
        \begin{itemize}
    \else
       \begin{tcolorbox}[
    enhanced,
    breakable,
    colback=fvdcyanlight,
    colframe=fvdcyan,
    coltext=fvdtext,
    boxrule=0.5pt,
    leftrule=4pt,
    arc=4mm,
    outer arc=4mm,
    left=7mm,
    right=6mm,
    top=5mm,
    bottom=5mm,
    before skip=6mm,
    after skip=6mm,
    title=Voorkennis,
    coltitle=fvdcyan,
    fonttitle=\bfseries\Large,
    attach title to upper={\par\smallskip},
    boxed title style={
        empty,
        left=0mm,
        right=0mm,
        top=0mm,
        bottom=0mm
    },
    overlay={
        \draw[
            fvdcyan,
            line width=0.5pt
        ]
        ([xshift=42mm,yshift=-13mm]frame.north west)
        --
        ([xshift=-7mm,yshift=-13mm]frame.north east);
    }
]
        \begin{itemize}
    \fi
}
{%
    \end{itemize}
    \ifdefined\HCode
        \HCode{%
</div>
</section>
        }%
    \else
        \end{tcolorbox}
    \fi
}


\newenvironment{leerdoelen}
{%
    \ifdefined\HCode
        \HCode{%
<section class="learning-box learning-box--goals">
<div class="learning-box__header">
<span class="learning-box__icon" aria-hidden="true"></span>
<span class="learning-box__title">Leerdoelen</span>
</div>
<div class="learning-box__content">
        }%
        \begin{itemize}
    \else
\begin{tcolorbox}[
    enhanced,
    breakable,
    colback=fvdmagentalight,
    colframe=fvdmagenta,
    coltext=fvdtext,
    boxrule=0.5pt,
    leftrule=4pt,
    arc=4mm,
    outer arc=4mm,
    left=7mm,
    right=6mm,
    top=5mm,
    bottom=5mm,
    before skip=6mm,
    after skip=6mm,
    title=Leerdoelen,
    coltitle=fvdmagenta,
    fonttitle=\bfseries\Large,
    attach title to upper={\par\smallskip},
    boxed title style={
        empty,
        left=0mm,
        right=0mm,
        top=0mm,
        bottom=0mm
    },
    overlay={
        \draw[
            fvdmagenta,
            line width=0.5pt
        ]
        ([xshift=42mm,yshift=-13mm]frame.north west)
        --
        ([xshift=-7mm,yshift=-13mm]frame.north east);
    }
]
        \begin{itemize}
    \fi
}
{%
    \end{itemize}
    \ifdefined\HCode
        \HCode{%
</div>
</section>
        }%
    \else
        \end{tcolorbox}
    \fi
}


\addPrintStyle{../../..}

\title{Oefeningen — Periodieke functies en amplitude}
\author{Ingmar Herreman}

\begin{document}

% FVD_CHAPTER_DATA_START
% Automatisch gegenereerd uit index.tex.
% Niet handmatig aanpassen.
\fvdchapterdata{5}{Het gedrag van functies beschrijven}{5}
% FVD_CHAPTER_DATA_END




\fvdchapterdata{5}{Het gedrag van functies beschrijven}{5}

\maketitle
\FVDbeginExercises

\begin{opmerking}
De oefeningen met \(\blacklozenge\) zijn essentieel. Het leerplandoel vraagt
\emph{analyseren}: je moet periode en amplitude niet alleen berekenen, maar ook
afleiden, voorstellingen verbinden en in context interpreteren.
\end{opmerking}

\FVDsection{Basis — periode en amplitude bepalen}

\begin{oefening}[Herhaling herkennen]{\oefB\ESS}
Een functie heeft maxima bij
\[
x=-5,\,-1,\,3,\,7,\,11
\]
en de tussenliggende grafiekdelen herhalen telkens hetzelfde patroon.

\begin{question}
Bepaal de periode.
\begin{oplossing}
De afstand tussen opeenvolgende overeenkomstige maxima is telkens \(4\).
Dus
\[
T=4.
\]
\end{oplossing}
\end{question}

\begin{question}
Is \(8\) ook een periode?
\begin{oplossing}
Ja. Na twee volledige periodes is de grafiek opnieuw dezelfde.
De fundamentele periode is echter \(4\).
\end{oplossing}
\end{question}
\end{oefening}

\begin{oefening}[Maximum, minimum en amplitude]{\oefB\ESS}
Een periodieke functie heeft
\[
M=13,\qquad m=5.
\]

\begin{question}
Bepaal de amplitude.
\begin{oplossing}
\[
A=\frac{13-5}{2}=4.
\]
\end{oplossing}
\end{question}

\begin{question}
Bepaal de middenwaarde.
\begin{oplossing}
\[
d=\frac{13+5}{2}=9.
\]
De middenlijn is \(y=9\).
\end{oplossing}
\end{question}
\end{oefening}

\begin{oefening}[Functiewaarden reconstrueren]{\oefB\ESS}
Een functie \(f\) heeft periode \(6\) en
\[
f(1)=3,\qquad f(4)=-2.
\]

Bepaal:
\[
f(7),\quad f(13),\quad f(-5),\quad f(10).
\]

\begin{oplossing}
Omdat verschuiven over \(6\) de functiewaarde niet verandert:
\[
f(7)=3,\qquad f(13)=3,\qquad f(-5)=3,\qquad f(10)=f(4)=-2.
\]
\end{oplossing}
\end{oefening}

\begin{oefening}[Horizontaal of verticaal?]{\oefB\ESS}
Vul aan met \emph{periode} of \emph{amplitude}.

\begin{question}
De horizontale lengte van één volledige herhaling.
\begin{oplossing}
Periode.
\end{oplossing}
\end{question}

\begin{question}
De maximale verticale uitwijking ten opzichte van de middenlijn.
\begin{oplossing}
Amplitude.
\end{oplossing}
\end{question}

\begin{question}
Eenheid bij een tijdgrafiek: seconden.
\begin{oplossing}
Dat kan de eenheid van de periode zijn.
\end{oplossing}
\end{question}
\end{oefening}

\FVDsection{Verdieping — analyseren en interpreteren}

\begin{oefening}[Is dit werkelijk één periode?]{\oefU\ESS}
Een leerling ziet twee opeenvolgende nulpunten van een golvende grafiek op
\(x=2\) en \(x=5\) en besluit:
\[
T=3.
\]

Leg uit waarom deze conclusie zonder extra informatie niet gerechtvaardigd is.

\begin{oplossing}
Opeenvolgende nulpunten hoeven geen overeenkomstige punten in het patroon te zijn.
De grafiek kan bij het ene nulpunt stijgend en bij het andere dalend door de as gaan.
Voor de periode moeten we punten vergelijken die dezelfde plaats en richting
binnen het herhalende patroon hebben, bijvoorbeeld twee opeenvolgende maxima.
\end{oplossing}
\end{oefening}

\begin{oefening}[Van één cyclus naar de grafiek]{\oefU\ESS}
Een periodieke functie heeft periode \(4\). Op \(0\leq x\leq4\) geldt:
\[
f(0)=1,\quad f(1)=4,\quad f(2)=1,\quad f(3)=-2,\quad f(4)=1.
\]

\begin{question}
Bepaal de amplitude en middenwaarde.
\begin{oplossing}
\[
M=4,\qquad m=-2.
\]
Dus
\[
A=\frac{4-(-2)}2=3,
\qquad
d=\frac{4+(-2)}2=1.
\]
\end{oplossing}
\end{question}

\begin{question}
Bepaal \(f(5)\), \(f(6)\), \(f(7)\) en \(f(8)\).
\begin{oplossing}
Door de periode \(4\):
\[
f(5)=4,\quad f(6)=1,\quad f(7)=-2,\quad f(8)=1.
\]
\end{oplossing}
\end{question}

\begin{question}
Schets minstens drie periodes van een mogelijke vloeiende grafiek.
\begin{oplossing}
De schets moet het gegeven patroon om de \(4\) eenheden exact herhalen.
Meerdere vloeiende grafieken zijn mogelijk omdat slechts enkele punten zijn gegeven.
\end{oplossing}
\end{question}
\end{oefening}

\begin{oefening}[Sensorsignaal]{\oefU\ESS}
Een periodiek sensorsignaal varieert tussen \(0{,}8\) V en \(5{,}2\) V.
Eenzelfde piek keert terug na \(0{,}040\) s.

\begin{question}
Bepaal de periode.
\begin{oplossing}
\[
T=0{,}040\ \mathrm{s}.
\]
\end{oplossing}
\end{question}

\begin{question}
Bepaal amplitude en middenwaarde.
\begin{oplossing}
\[
A=\frac{5{,}2-0{,}8}{2}=2{,}2\ \mathrm V,
\]
en
\[
d=\frac{5{,}2+0{,}8}{2}=3{,}0\ \mathrm V.
\]
\end{oplossing}
\end{question}

\begin{question}
Interpreteer beide kenmerken.
\begin{oplossing}
Het signaalpatroon herhaalt zich om de \(0{,}040\) s.
De spanning wijkt maximaal \(2{,}2\) V af van de middenwaarde \(3{,}0\) V.
\end{oplossing}
\end{question}
\end{oefening}

\begin{oefening}[Draaiend wiel]{\oefU\ESS}
Een punt op een wiel bereikt een maximale hoogte van \(2{,}6\) m en een minimale
hoogte van \(1{,}0\) m. Het wiel maakt één volledige omwenteling in \(3\) s.

\begin{question}
Bepaal de amplitude van de hoogtefunctie.
\begin{oplossing}
\[
A=\frac{2{,}6-1{,}0}{2}=0{,}8\ \mathrm m.
\]
\end{oplossing}
\end{question}

\begin{question}
Bepaal de middenwaarde en interpreteer ze.
\begin{oplossing}
\[
d=\frac{2{,}6+1{,}0}{2}=1{,}8\ \mathrm m.
\]
Dit is de hoogte van het middelpunt van het wiel.
\end{oplossing}
\end{question}

\begin{question}
Bepaal de periode en interpreteer ze.
\begin{oplossing}
\[
T=3\ \mathrm s.
\]
Na \(3\) seconden heeft het punt opnieuw dezelfde plaats in de cyclus.
\end{oplossing}
\end{question}
\end{oefening}

\begin{oefening}[Foutenanalyse]{\oefU\ESS}
Een grafiek oscilleert tussen \(y=-3\) en \(y=7\).
Een leerling zegt:
``De amplitude is \(10\), want \(7-(-3)=10\).''

Analyseer de fout.

\begin{oplossing}
\(10\) is de totale verticale afstand tussen minimum en maximum.
De amplitude is de helft daarvan:
\[
A=\frac{7-(-3)}2=5.
\]
De middenwaarde is
\[
d=\frac{7+(-3)}2=2.
\]
\end{oplossing}
\end{oefening}

\begin{oefening}[Zelfde amplitude, andere periode]{\oefU\ESS}
Twee periodieke functies hebben beide maximum \(5\) en minimum \(-1\).
De eerste heeft periode \(4\), de tweede periode \(10\).

\begin{question}
Vergelijk hun amplitudes.
\begin{oplossing}
Beide hebben
\[
A=\frac{5-(-1)}2=3.
\]
\end{oplossing}
\end{question}

\begin{question}
Kun je besluiten dat hun grafieken gelijk zijn?
\begin{oplossing}
Nee. Ze hebben al een verschillende periode. Bovendien bepalen periode,
maximum en minimum de precieze vorm van een periodieke grafiek niet volledig.
\end{oplossing}
\end{question}
\end{oefening}

\FVDsection{Gevorderd — onderzoeken en redeneren}

\begin{oefening}[Hoeveel informatie is genoeg?]{\oefG}
Van een periodieke functie weet je:
\[
T=6,\qquad A=2,\qquad d=1.
\]

Is daarmee de grafiek uniek bepaald? Motiveer en teken twee verschillende
mogelijke grafieken die aan de gegevens voldoen.

\begin{oplossing}
Nee. De gegevens bepalen de horizontale herhalingslengte en de verticale schaal,
maar niet de vorm binnen één periode. Er zijn dus meerdere periodieke grafieken
mogelijk met dezelfde \(T\), \(A\) en \(d\).
\end{oplossing}
\end{oefening}

\begin{oefening}[Een verborgen kleinere periode]{\oefG}
Een leerling toont aan dat
\[
f(x+12)=f(x)
\]
voor alle \(x\) en concludeert dat de fundamentele periode \(12\) is.

Welke extra vraag moet nog onderzocht worden?

\begin{oplossing}
Er moet worden nagegaan of er een kleinere positieve waarde \(T<12\) bestaat
waarvoor eveneens
\[
f(x+T)=f(x)
\]
voor alle toegelaten \(x\). De gegeven gelijkheid toont alleen dat \(12\)
een periode is, niet noodzakelijk de kleinste positieve periode.
\end{oplossing}
\end{oefening}

\begin{oefening}[GeoGebra: periode versus amplitude]{\oefG}
Onderzoek:
\[
f(x)=\cos x,\qquad
g(x)=3\cos x,\qquad
h(x)=\cos(2x).
\]

\begin{enumerate}
  \item Bepaal grafisch periode en amplitude.
  \item Vergelijk \(f\) met \(g\): welk kenmerk verandert?
  \item Vergelijk \(f\) met \(h\): welk kenmerk verandert?
  \item Formuleer een vermoeden over de rol van de factor vóór de functie
        en de factor bij \(x\).
\end{enumerate}

\begin{oplossing}
\(f\) heeft amplitude \(1\) en periode \(2\pi\).
\(g\) heeft amplitude \(3\) en dezelfde periode \(2\pi\).
\(h\) heeft amplitude \(1\) en periode \(\pi\).

De factor vóór de cosinus verandert hier de verticale uitwijking.
De factor bij \(x\) verandert hier de horizontale herhaling.
De algemene transformatietheorie wordt later systematisch behandeld.
\end{oplossing}
\end{oefening}



\end{document}
